\documentclass{article}
\usepackage{amsmath, amssymb}

\title{VGASP-Snowman: Mathematical Backbone}
\author{Borealis S. Hedling}
\date{06 October 2026}

\begin{document}
\maketitle

\section{Introduction}
The VGASP-Snowman model provides a playful yet rigorous demonstration of
deterministic AI steering. The Snowman metaphor encodes governance concepts
such as invariants, operator ordering, dissipation, stability, constraint
fields, and multi-module linkage. This document formalizes these structures
mathematically.

\section{Invariant Manifold}
The Snowman is represented as a three-sphere invariant manifold:
\[
I(x) = \{ x \in \mathbb{R}^3 \mid r_1, r_2, r_3 > 0 \},
\]
where each radius \( r_i \) corresponds to a structural invariant. These
invariants define the allowable configuration space for generative behavior.

\section{Operator Ordering}
Snowman assembly follows a strict operator ordering:
\[
O_1 \prec O_2 \prec O_3,
\]
with:
\[
O_1 = \text{roll}, \quad
O_2 = \text{stack}, \quad
O_3 = \text{decorate}.
\]
This ordering enforces deterministic sequencing and prevents generative drift.

\section{Constraint Field}
A constraint field governs permissible transformations:
\[
C(x) \in \mathbb{R}^n,
\]
where each constraint \( C_i(x) \) corresponds to a governance rule. The field
ensures that generative outputs remain within safe and bounded domains.

\section{Dissipation Model}
Melting is modeled as a dissipation process:
\[
\frac{dS}{dt} \le 0,
\]
where \( S \) denotes structural stability. Dissipation ensures that destructive
transformations follow predictable thermodynamic behavior.

\section{Stability Condition}
Snowman stability requires:
\[
\rho(J) < 1,
\]
where \( J \) is the Jacobian of the generative transformation. This condition
ensures local stability and prevents runaway generative cascades.

\section{Module Linkage}
Arms, nose, and decorations attach via module linkage:
\[
M_i \leftrightarrow M_j,
\]
ensuring multi-module consistency. Linkage rules prevent incompatible module
combinations and enforce structural coherence.

\section{Banned Constructs}
Forbidden generative constructs are encoded as:
\[
B = \{ b \mid b \notin \text{SafeDomain} \},
\]
and enforced via VGASP’s governance firewall. This prevents unsafe or
non-compliant generative outputs.

\section{Conclusion}
The VGASP-Snowman mathematical backbone demonstrates how playful metaphors can
encode rigorous governance principles. Invariants, operator ordering,
constraint fields, dissipation, stability, and module linkage collectively
illustrate deterministic AI steering in a system-agnostic manner.

\section{References}

\begin{thebibliography}{9}

\bibitem{dai2026assured}
T.~Dai, D.~Simchi-Levi, M.~X.~Wu, and Y.~Xie,
\textit{Assured autonomy: How operations research powers and orchestrates generative AI systems}.
Production and Operations Management, vol.~35, no.~10, pp.~3579–3596, 2026.

\bibitem{smith2026deterministic}
C.~B.~Smith and D.~McCarthy,
\textit{Deterministic governance for generative systems: a policy-aligned runtime for validating AI outputs}.
AI and Ethics, Springer Nature, 2026.

\bibitem{kuai2026governance}
J.~Kuai and F.~Ferrari,
\textit{Generative AI governance: from regulatory design to institutional reordering}.
Information, Communication \& Society, vol.~29, no.~6, pp.~1803–1812, 2026.

\bibitem{dai2025arxiv}
T.~Dai, D.~Simchi-Levi, M.~X.~Wu, and Y.~Xie,
\textit{Assured Autonomy: How Operations Research Powers and Orchestrates Generative AI Systems}.
arXiv preprint arXiv:2512.23978, 2025.

\end{thebibliography}

\end{document}
